\chapter{روش نیوتن، الگوریتم و ارزیابی عددی}

\section{استخراج فرمول نیوتن}
بر اساس رابطه \cref{eq:h-formula}، اگر $f'(x_k)\neq 0$ باشد، تکرار نیوتن به صورت زیر تعریف می‌شود:
\begin{equation}
 x_{k+1}=x_k-\frac{f(x_k)}{f'(x_k)}.
 \label{eq:newton-main}
\end{equation}
از دید هندسی، ابتدا در نقطه $(x_k,f(x_k))$ خط مماس رسم می‌شود. سپس محل برخورد این خط با محور افقی به عنوان $x_{k+1}$ انتخاب می‌شود. معادله خط مماس چنین است:
\begin{equation}
 y=f(x_k)+f'(x_k)(x-x_k).
 \label{eq:tangent-line}
\end{equation}
با قرار دادن $y=0$ در \cref{eq:tangent-line} دوباره همان فرمول \cref{eq:newton-main} به دست می‌آید.

\section{شبه‌کد الگوریتم}
برای اجرای عددی روش نیوتن، باید مقدار اولیه، تابع، مشتق تابع، حد تحمل و حداکثر تعداد تکرار مشخص شود. شبه‌کد \cref{alg:newton} مراحل کلی الگوریتم را نشان می‌دهد.

\begin{algorithm}[H]
\caption{شبه‌کد روش نیوتن برای حل $f(x)=0$}
\label{alg:newton}
\begin{latin}
\begin{tabular}{ll}
\textbf{Input:} & function $f$, derivative $f'$, initial point $x_0$, tolerance $\varepsilon$, maximum iteration $N$ \\
\textbf{Output:} & approximate root $x_{k+1}$ \\
\hline
1. & for $k=0,1,\ldots,N-1$ do \\
2. & \quad if $f'(x_k)=0$, stop and report failure \\
3. & \quad $x_{k+1}=x_k-\dfrac{f(x_k)}{f'(x_k)}$ \\
4. & \quad if $|x_{k+1}-x_k|<\varepsilon$ or $|f(x_{k+1})|<\varepsilon$, return $x_{k+1}$ \\
5. & \quad set $x_k\leftarrow x_{k+1}$ \\
6. & end for \\
7. & report that the method did not converge within $N$ iterations.
\end{tabular}
\end{latin}
\end{algorithm}

\section{فلوچارت اجرای روش}
فلوچارت \cref{fig:flowchart} مسیر تصمیم‌گیری الگوریتم را نشان می‌دهد. این شکل به ویژه برای پیاده‌سازی برنامه مفید است، چون هم حالت موفق و هم حالت شکست را جدا می‌کند.

\begin{figure}[H]
\centering
\resizebox{0.72\textwidth}{!}{%
\begin{tikzpicture}[
node distance=1.05cm,
font=\small,
startstop/.style={ellipse, draw, minimum width=2.5cm, minimum height=0.65cm, align=center},
process/.style={rectangle, draw, rounded corners, minimum width=3.6cm, minimum height=0.65cm, align=center},
decision/.style={diamond, draw, aspect=2.15, inner sep=1pt, minimum width=3.5cm, align=center},
arrow/.style={-{Stealth[length=2.1mm]}, thick}
]
\node[startstop] (start) {شروع};
\node[process, below=of start] (input) {دریافت $f$، $f'$، $x_0$ و $\varepsilon$};
\node[decision, below=of input] (deriv) {آیا $f'(x_k)=0$؟};
\node[process, below=of deriv] (update) {محاسبه $x_{k+1}=x_k-\frac{f(x_k)}{f'(x_k)}$};
\node[decision, below=of update] (stop) {آیا شرط توقف برقرار است؟};
\node[startstop, below=of stop] (success) {پایان: ریشه تقریبی};
\node[startstop, right=2.55cm of deriv] (fail) {شکست};
\node[process, left=2.55cm of stop] (next) {$k\leftarrow k+1$};
\draw[arrow] (start) -- (input);
\draw[arrow] (input) -- (deriv);
\draw[arrow] (deriv) -- node[right]{خیر} (update);
\draw[arrow] (deriv) -- node[above]{بله} (fail);
\draw[arrow] (update) -- (stop);
\draw[arrow] (stop) -- node[right]{بله} (success);
\draw[arrow] (stop) -- node[above]{خیر} (next);
\draw[arrow] (next) |- (deriv);
\end{tikzpicture}%
}
\caption{فلوچارت ساده اجرای روش نیوتن}
\label{fig:flowchart}
\end{figure}

\section{قضیه همگرایی درجه دوم}
در این بخش نشان می‌دهیم که اگر نقطه شروع به ریشه ساده نزدیک باشد، خطای روش نیوتن تقریباً با مربع خطای قبلی متناسب است. این همان دلیل سرعت زیاد روش است.

\begin{theorem}[همگرایی درجه دوم روش نیوتن]
\label{thm:quadratic}
فرض کنید $f$ در همسایگی ریشه ساده $\alpha$ دوبار پیوسته مشتق‌پذیر باشد، $f(\alpha)=0$ و $f'(\alpha)\neq 0$. اگر $x_0$ به اندازه کافی به $\alpha$ نزدیک باشد، آنگاه دنباله حاصل از \cref{eq:newton-main} به $\alpha$ همگرا می‌شود و برای خطا داریم
\begin{equation}
 |e_{k+1}|\leq C|e_k|^2,
 \label{eq:quadratic-bound}
\end{equation}
که در آن $C$ یک ثابت مثبت است.
\end{theorem}

\begin{proof}
با توجه به $e_k=x_k-\alpha$، تکرار نیوتن را از ریشه دقیق کم می‌کنیم. با استفاده از بسط تیلور حول $\alpha$ داریم
\begin{align}
 e_{k+1}
 &=x_{k+1}-\alpha \notag\\
 &=x_k-\frac{f(x_k)}{f'(x_k)}-\alpha \notag\\
 &=e_k-\frac{f(\alpha)+f'(\alpha)e_k+\frac{1}{2}f''(\xi_k)e_k^2}{f'(x_k)} \notag\\
 &=e_k-\frac{f'(\alpha)e_k}{f'(x_k)}-
 \frac{f''(\xi_k)}{2f'(x_k)}e_k^2 .
 \label{eq:error-expansion-long}
\end{align}
از طرفی چون $f'$ پیوسته است و $f'(\alpha)\neq 0$، برای $x_k$های نزدیک به $\alpha$، مقدار $f'(x_k)$ از صفر دور می‌ماند. بنابراین می‌توان ثابت مثبتی مانند $C$ پیدا کرد که جمله‌های سمت راست \cref{eq:error-expansion-long} را کنترل کند و رابطه \cref{eq:quadratic-bound} نتیجه شود. پس مرتبه همگرایی روش نیوتن در نزدیکی ریشه ساده برابر دو است.
\end{proof}

فرمول چندخطی \cref{eq:error-expansion-long} نمونه‌ای از یک رابطه طولانی است که مطابق اصول نگارش ریاضی، در چند خط شکسته و شماره‌گذاری شده است.

\section{مثال عددی و جدول تکرارها}
تابع زیر را در نظر بگیرید:
\begin{equation}
 f(x)=x^3-x-2,
 \qquad
 f'(x)=3x^2-1.
 \label{eq:example-function}
\end{equation}
برای مقدار اولیه $x_0=1.5$، روش نیوتن طبق \cref{eq:newton-main} به ریشه حقیقی تابع نزدیک می‌شود. مقدار دقیق ریشه تقریباً برابر است با
\begin{equation}
 \alpha\approx 1.5213797068.
 \label{eq:alpha-value}
\end{equation}
نتیجه چند تکرار اول در \cref{tab:newton-table} آمده است.

\begin{table}[H]
\centering
\caption{جدول تکرارهای روش نیوتن برای تابع $f(x)=x^3-x-2$}
\label{tab:newton-table}
\begin{tabular}{|c|c|c|c|c|}
\hline
$k$ & $x_k$ & $f(x_k)$ & $x_{k+1}$ & $|x_k-\alpha|$ \\
\hline
$0$ & $1.500000$ & $-0.125000$ & $1.521739$ & $2.1379\times10^{-2}$ \\
\hline
$1$ & $1.521739$ & $2.1369\times10^{-3}$ & $1.521380$ & $3.5942\times10^{-4}$ \\
\hline
$2$ & $1.521380$ & $5.8939\times10^{-7}$ & $1.521380$ & $9.9160\times10^{-8}$ \\
\hline
$3$ & $1.521380$ & $4.5297\times10^{-14}$ & $1.521380$ & $7.5495\times10^{-15}$ \\
\hline
$4$ & $1.521380$ & $0$ & $1.521380$ & $0$ \\
\hline
\end{tabular}
\end{table}

همان‌طور که در \cref{tab:newton-table} دیده می‌شود، خطا پس از چند گام به سرعت کوچک می‌شود. این رفتار با نتیجه نظری \cref{thm:quadratic} سازگار است.

\section{شکل‌های کمکی و تفسیر هندسی}
در \cref{fig:subfigures} دو تصویر کمکی آورده شده است. تصویر سمت راست ایده هندسی خط مماس را نشان می‌دهد و تصویر سمت چپ کاهش سریع خطا را به صورت شماتیک نمایش می‌دهد.

\begin{figure}[H]
\centering
\begin{subfigure}{0.46\textwidth}
\centering
\begin{tikzpicture}[scale=1.0]
\draw[->] (-0.2,0) -- (4.2,0) node[right] {$x$};
\draw[->] (0,-0.2) -- (0,3.2) node[above] {$|e_k|$};
\draw (0.5,2.8) node[circle,fill,inner sep=1.5pt]{};
\draw (1.4,1.2) node[circle,fill,inner sep=1.5pt]{};
\draw (2.3,0.35) node[circle,fill,inner sep=1.5pt]{};
\draw (3.2,0.08) node[circle,fill,inner sep=1.5pt]{};
\draw[dashed] (0.5,2.8) -- (1.4,1.2) -- (2.3,0.35) -- (3.2,0.08);
\node at (2.3,2.7) {کاهش خطا};
\end{tikzpicture}
\caption{رفتار شماتیک خطا}
\label{fig:error-sub}
\end{subfigure}
\hfill
\begin{subfigure}{0.46\textwidth}
\centering
\begin{tikzpicture}[scale=1.0]
\draw[->] (-0.2,0) -- (4.3,0) node[right] {$x$};
\draw[->] (0,-1.1) -- (0,2.4) node[above] {$y$};
\draw[domain=0.7:3.4,smooth,variable=\x] plot ({\x},{0.45*((\x-2.35)^3+0.6*(\x-2.35))});
\draw[dashed] (2.0,-0.35) -- (2.0,0.35);
\draw (2.0,0.35) node[circle,fill,inner sep=1.4pt]{} node[above right] {$x_k$};
\draw (2.0,0.35) -- (2.55,0);
\draw (2.55,0) node[circle,fill,inner sep=1.4pt]{} node[below] {$x_{k+1}$};
\node at (2.8,1.55) {مماس};
\end{tikzpicture}
\caption{تقاطع مماس با محور افقی}
\label{fig:tangent-sub}
\end{subfigure}
\caption{دو زیرشکل برای توضیح هندسی و عددی روش نیوتن}
\label{fig:subfigures}
\end{figure}

\section{جمع‌بندی}
روش نیوتن نمونه‌ای روشن از پیوند میان تحلیل ریاضی و الگوریتم محاسباتی است. از یک طرف، فرمول آن از بسط تیلور و خط مماس به دست می‌آید؛ از طرف دیگر، اجرای آن فقط به یک رابطه تکراری ساده نیاز دارد. بر اساس \cref{thm:quadratic}، اگر نقطه شروع مناسب انتخاب شود و ریشه ساده باشد، همگرایی روش درجه دوم خواهد بود. البته باید به شرط $f'(x_k)\neq 0$ و انتخاب مقدار اولیه نیز توجه کرد؛ زیرا انتخاب نامناسب می‌تواند باعث واگرایی یا توقف الگوریتم شود \cite{nocedal,rudin}.
